% =============================================================================
% CHƯƠNG 4: TRIỂN KHAI TRÊN MININET
% =============================================================================
\chapter{Triển khai trên Mininet}

\section{Môi trường cài đặt}

\subsection{Yêu cầu hệ thống}

\begin{table}[htbp]
    \centering
    \caption{Yêu cầu phần mềm và phần cứng}
    \vspace{0.2cm}
    \begin{tabularx}{\textwidth}{|>{\raggedright\arraybackslash}p{4cm}|>{\raggedright\arraybackslash}X|}
        \hline
        \textbf{Thành phần} & \textbf{Phiên bản / Yêu cầu} \\ \hline
        Hệ điều hành & Ubuntu 22.04 LTS (64-bit) \\ \hline
        RAM & Tối thiểu 4 GB (khuyến nghị 8 GB) \\ \hline
        CPU & 2 cores trở lên \\ \hline
        Disk & 20 GB trống \\ \hline
        Python & Python 3.10+ \\ \hline
        Mininet & 2.3.0+ \\ \hline
        Open vSwitch & 2.17+ \\ \hline
        iperf3 & 3.9+ \\ \hline
        iproute2 & Hỗ trợ MPLS (kernel 4.1+) \\ \hline
    \end{tabularx}
\end{table}

\subsection{Hướng dẫn cài đặt từng bước}

\textbf{Bước 1: Cập nhật hệ thống}

\begin{lstlisting}[language=bash, frame=single, basicstyle=\ttfamily\small]
sudo apt update && sudo apt upgrade -y
\end{lstlisting}

\textbf{Bước 2: Cài đặt các gói cần thiết}

\begin{lstlisting}[language=bash, frame=single, basicstyle=\ttfamily\small]
sudo apt install -y \
    mininet \
    openvswitch-switch openvswitch-common \
    python3 python3-pip \
    iperf3 iperf \
    tcpdump net-tools iproute2 \
    iputils-ping traceroute \
    python3-matplotlib python3-numpy \
    curl wget git
\end{lstlisting}

\textbf{Bước 3: Cài thư viện Python bổ sung}

\begin{lstlisting}[language=bash, frame=single, basicstyle=\ttfamily\small]
pip3 install matplotlib numpy scipy
\end{lstlisting}

\textbf{Bước 4: Bật Open vSwitch}

\begin{lstlisting}[language=bash, frame=single, basicstyle=\ttfamily\small]
sudo systemctl enable openvswitch-switch
sudo systemctl start openvswitch-switch
# Kiểm tra:
sudo ovs-vsctl show
\end{lstlisting}

\textbf{Bước 5: Kích hoạt MPLS kernel module}

\begin{lstlisting}[language=bash, frame=single, basicstyle=\ttfamily\small]
# Load module MPLS
sudo modprobe mpls_router
sudo modprobe mpls_iptunnel

# Kiểm tra module đã load:
lsmod | grep mpls

# Bật platform labels (cần thiết cho iproute2 MPLS)
echo 'net.mpls.platform_labels = 1000' | sudo tee -a /etc/sysctl.conf
sudo sysctl -p

# Xác nhận:
sysctl net.mpls.platform_labels
\end{lstlisting}

\textbf{Bước 6: Kiểm tra Mininet hoạt động}

\begin{lstlisting}[language=bash, frame=single, basicstyle=\ttfamily\small]
# Test đơn giản: topo single,3 với ping
sudo mn --topo single,3 --test ping

# Kết quả mong đợi:
# *** Ping: testing ping reachability
# h1 -> h2 h3
# h2 -> h1 h3
# h3 -> h1 h2
# *** Results: 0% dropped (6/6 received)

# Dọn dẹp sau test
sudo mn -c
\end{lstlisting}

\section{Cấu trúc file code}

\begin{verbatim}
CuoiKy/
\u251c\u2500\u2500 source/                      <- Tat ca file Python Mininet
|   \u251c\u2500\u2500 mpls_metro_topology.py   <- Topology tong the (file chinh)
|   \u251c\u2500\u2500 flat.py                  <- Tham khao: Branch 1 Flat
|   \u251c\u2500\u2500 coreditrubutionaccess.py <- Tham khao: Branch 2 3-Tier
|   \u2514\u2500\u2500 spineleaf.py             <- Tham khao: Branch 3 Leaf-Spine
\u251c\u2500\u2500 scripts/
|   \u251c\u2500\u2500 run_tests.sh             <- Script 3 kich ban do luong
|   \u251c\u2500\u2500 collect_results.py       <- Phan tich log, ve bieu do
|   \u2514\u2500\u2500 mpls_config.sh           <- Cau hinh MPLS ngoai Mininet
\u251c\u2500\u2500 results/
|   \u251c\u2500\u2500 raw/                     <- Log iperf3 va ping
|   \u251c\u2500\u2500 csv/                     <- So lieu tong hop CSV
|   \u2514\u2500\u2500 charts/                  <- Bieu do PNG
\u251c\u2500\u2500 source_latex/                <- Bao cao LaTeX
\u251c\u2500\u2500 docs/
|   \u251c\u2500\u2500 debai.txt
|   \u251c\u2500\u2500 kienthuc.txt
|   \u2514\u2500\u2500 huong_dan_cai_dat.md
\u2514\u2500\u2500 README.md
\end{verbatim}

\section{Hướng dẫn chạy topology}

\subsection{Khởi động topology tổng thể}

\begin{lstlisting}[language=bash, frame=single, basicstyle=\ttfamily\small]
# Vao thu muc source (chua tat ca file .py)
cd ~/CuoiKy/source

# Don dep Mininet cu (neu can)
sudo mn -c

# Chay topology tong the (3 chi nhanh + MPLS backbone)
sudo python3 mpls_metro_topology.py

# === Hoac chay tung chi nhanh doc lap ===
sudo python3 flat.py              # Branch 1 - Flat
sudo python3 coreditrubutionaccess.py  # Branch 2 - 3-Tier
sudo python3 spineleaf.py         # Branch 3 - Leaf-Spine
\end{lstlisting}

\subsection{Kiểm tra kết nối trong Mininet CLI}

Sau khi topology khởi động, bạn sẽ thấy dấu nhắc \texttt{mininet>}:

\begin{lstlisting}[language=bash, frame=single, basicstyle=\ttfamily\small]
# Kiểm tra tất cả node
mininet> nodes

# Kiểm tra toàn bộ liên kết
mininet> net

# Ping từ chi nhánh 1 sang chi nhánh 2
mininet> h1a ping -c 5 10.2.0.1

# Ping từ chi nhánh 1 sang chi nhánh 3
mininet> h1a ping -c 5 10.3.0.1

# Ping tất cả (có thể thời gian dài)
mininet> pingall

# Xem bảng định tuyến của CE1
mininet> ce1 ip route

# Xem bảng định tuyến của P1 (MPLS core)
mininet> p1 ip route

# Kiểm tra MPLS labels (nếu kernel hỗ trợ)
mininet> p1 ip -f mpls route

# Traceroute từ h1a → h2a (xem đường qua backbone)
mininet> h1a traceroute -n 10.2.0.1
\end{lstlisting}

\subsection{Đo lường iperf3 trong CLI}

\begin{lstlisting}[language=bash, frame=single, basicstyle=\ttfamily\small]
# Đo TCP throughput: h1a → h2a
mininet> h2a iperf3 -s &
mininet> h1a iperf3 -c 10.2.0.1 -t 10

# Đo UDP jitter và packet loss: h1a → h3a
mininet> h3a iperf3 -s &
mininet> h1a iperf3 -c 10.3.0.1 -u -b 50M -t 10

# Lưu kết quả ra file (cần chạy trong terminal riêng)
# ip netns exec h1a iperf3 -c 10.2.0.1 -t 10 -J > results.json
\end{lstlisting}

\subsection{Thoát và dọn dẹp}

\begin{lstlisting}[language=bash, frame=single, basicstyle=\ttfamily\small]
# Trong Mininet CLI:
mininet> exit

# Hoặc Ctrl+D

# Dọn dẹp toàn bộ:
sudo mn -c
\end{lstlisting}

\section{Mô phỏng MPLS Label Switching}

\subsection{Cơ chế mô phỏng}

Do giới hạn của Mininet (không có phần cứng MPLS thật), đề tài mô phỏng MPLS theo hai cách:

\textbf{Cách 1: Linux Kernel MPLS (iproute2)} – Dùng khi kernel \texttt{\textgreater{}= 4.1} hỗ trợ module \texttt{mpls\_router}:

\begin{lstlisting}[language=bash, frame=single, basicstyle=\ttfamily\small]
# Bật MPLS trên interface
sysctl -w net.mpls.platform_labels=1000
sysctl -w net.mpls.conf.eth0.input=1

# PUSH: PE1 gắn label 100 cho traffic đến 10.2.0.0/24
ip route add 10.2.0.0/24 encap mpls 100/200 \
    via inet 10.0.12.2 dev pe1-p1

# SWAP: P1 swap label 100 → 200, chuyển sang P2
ip -f mpls route add 100 as 200 via inet 10.0.99.2 dev p1-p2

# POP: PE2 pop label, giao về CE2
ip -f mpls route add 200 via inet 10.100.2.1 dev pe2-ce2
\end{lstlisting}

\textbf{Cách 2: Static Routing (IP Forwarding)} – Dùng như fallback khi kernel MPLS không khả dụng. Định tuyến tĩnh đảm bảo gói tin đi đúng đường như MPLS LSP mô tả.

\subsection{Xác minh hoạt động MPLS}

\begin{lstlisting}[language=bash, frame=single, basicstyle=\ttfamily\small]
# Xem MPLS forwarding table của P1
ip netns exec p1 ip -f mpls route

# Bắt gói tin trên interface P1 để xem label
ip netns exec p1 tcpdump -i p1-pe1 -n -v mpls

# Kiểm tra đường đi gói tin qua backbone
ip netns exec h1a traceroute -n 10.2.0.1
\end{lstlisting}
